% Generated draft: requires mathematical and source review.
\documentclass[11pt]{article}
\usepackage[margin=25mm]{geometry}
\usepackage{fontspec}
\setmainfont{Latin Modern Roman}
\usepackage{amsmath,amssymb,mathtools,unicode-math}
\setmathfont{Latin Modern Math}
\usepackage{enumitem,needspace,xurl,hyperref}
\setlength{\parindent}{0pt}
\setlength{\parskip}{5pt}
\setlength{\emergencystretch}{3em}
\newcommand{\N}{\mathbb{N}}
\newcommand{\Z}{\mathbb{Z}}
\newcommand{\Q}{\mathbb{Q}}
\newcommand{\R}{\mathbb{R}}
\newcommand{\C}{\mathbb{C}}
\newcommand{\F}{\mathbb{F}}
\newcommand{\A}{\mathbb{A}}
\newcommand{\PP}{\mathbb{P}}
\newcommand{\E}{\mathbb{E}}
\newcommand{\eps}{\varepsilon}
\newcommand{\dd}{\,\mathrm{d}}
\begin{document}
\section*{UnsolvedMath ID: 30002411}
\textbf{Source problem code:} catalogue entry
\subsection{Definitions and mathematical statement}
An involution constraint for an evolution system is a differential constraint on its unknown fields that is propagated by the evolution for sufficiently regular solutions satisfying the relevant compatibility conditions. Schematically, writing the fields as \(U\) and the constraint as \(\mathcal C(U)=0\), propagation means that satisfaction of the constraint initially implies its satisfaction at subsequent times for which the solution and the required hypotheses hold. This notation describes the terminology only; the supplied record specifies neither an evolution operator nor a constraint operator.

\textbf{Source question.} Can high-resolution numerical methodology for involution-constrained systems be applied to equations coupled with Maxwell's equations, such as elastodynamics?

The intended research task is to investigate the extension of such numerical methodology to Maxwell-coupled systems, with elastodynamics named as an example. The source does not specify the coupled equations, constitutive relations, spatial domain, initial or boundary conditions, solution regularity, numerical method, or a mathematical criterion for successful application. In particular, neither the meaning of high resolution nor the required form of constraint preservation is fixed. Consequently, this is a research question rather than a fully quantified existence, convergence, stability, or accuracy assertion; a unique precise theorem cannot be extracted without additional hypotheses and a specified target property.
\subsection{Short English statement}
The question asks whether high-resolution methods for evolution equations with propagated differential constraints can be extended to systems coupled with Maxwell's equations, including elastodynamics. The particular coupling and the properties the numerical method must satisfy remain unspecified. The supplied partially solved label is a dated curation assessment, not an independently verified mathematical conclusion.
\subsection{Sources}
\begin{itemize}
\item Local UnsolvedMath registry, permanent ID \(30002411\), source problem code \texttt{catalogue entry}, title ``Constraint-Preserving Methods for Maxwell-Coupled Systems,'' in the Oberwolfach Reports Open Problems collection. No individual proposer is identified in the supplied record.
\item Supplied report citation: \emph{High-Resolution Mathematical and Numerical Analysis of Involution-Constrained PDEs} (2014), Oberwolfach Reports, pp.\ 2691--2747, DOI \texttt{10.4171/owr/2013/47}. \url{https://doi.org/10.4171/owr/2013/47}. The registry records proposed year 2013; the supplied citation gives 2014. The background describes Oberwolfach Reports as documenting mathematical workshops and their open-problem discussions.
\item Supplied literature source, ``Oberwolfach Report'': \url{https://ems.press/content/serial-article-files/46478}.
\item Supplied literature source, ``Dumbser--Chiocchetti--Peshkov'': \url{https://arxiv.org/abs/2003.02799}. No title or more detailed bibliographic attribution is supplied for this reference.
\item Literature checked at \texttt{2026-08-21}: this is supplied dated curation, not independently verified here. Its status is ``partially solved.'' Its assessment states that differential-form and mimetic formulations provide a route for coupled elastodynamic systems, that curl-preserving schemes exist, and that no general methodology covering the extracted question is established. The background expressly describes this as dated literature triage derived from the supplied OWR curation and requests independent verification of time-sensitive status and cited claims.
\item The supplied registry marks the statement as exact and reports that manual review found the extracted wording self-contained, mathematically coherent, and free of LaTeX or report-layout artifacts. These are supplied verification claims; they do not resolve the missing hypotheses or establish mathematical correctness.
\end{itemize}
\end{document}
